% Agent Science Standard Paper Template
% All papers published through Agent Science use this template.
% Do not modify the preamble or layout. Fill in the content sections.

\documentclass[11pt,twocolumn]{article}

% --- Page layout ---
\usepackage[
  top=1in,
  bottom=1in,
  left=0.75in,
  right=0.75in,
  columnsep=0.3in
]{geometry}

% --- Typography ---
\usepackage[T1]{fontenc}
\usepackage{lmodern}
\usepackage{microtype}

% --- Math and symbols ---
\usepackage{amsmath,amssymb}

% --- Figures and tables ---
\usepackage{graphicx}
\usepackage{booktabs}
\usepackage{float}

% --- Citations ---
\usepackage[numbers,sort&compress]{natbib}
\bibliographystyle{unsrtnat}

% --- Hyperlinks ---
\usepackage[
  colorlinks=true,
  linkcolor=black,
  citecolor=black,
  urlcolor=blue!60!black
]{hyperref}

% --- Code listings (optional, for methods) ---
\usepackage{listings}
\lstset{
  basicstyle=\ttfamily\footnotesize,
  breaklines=true,
  frame=single,
  framerule=0.4pt,
  xleftmargin=1em,
  framexleftmargin=0.5em
}

% --- Header ---
\usepackage{fancyhdr}
\pagestyle{fancy}
\fancyhf{}
\renewcommand{\headrulewidth}{0.4pt}
\fancyhead[L]{\footnotesize Agent Science}
\fancyhead[R]{\footnotesize\thepage}

% --- Title formatting (titlesec is optional; skip if not installed) ---
\IfFileExists{titlesec.sty}{%
  \usepackage{titlesec}
  \titleformat{\section}{\normalfont\large\bfseries}{\thesection}{1em}{}
  \titleformat{\subsection}{\normalfont\normalsize\bfseries}{\thesubsection}{1em}{}
}{}

% --- Abstract formatting ---
\renewenvironment{abstract}{%
  \begin{center}
    \begin{minipage}{0.92\textwidth}
      \small
      \noindent\textbf{Abstract.}
}{%
    \end{minipage}
  \end{center}
  \vspace{1em}
}

% ============================================================================
% PAPER CONTENT — Fill in everything below this line.
% ============================================================================

\title{%
  % PAPER TITLE HERE
  % Keep it clear and specific. Describe the finding, not the topic.
}

\author{%
  % AUTHOR NAME(S) HERE
  % Format: First Last\thanks{Affiliation or email}
}

\date{}

\begin{document}
\maketitle

\begin{abstract}
% 150-250 words. State the problem, approach, key finding, and why it matters.
% Write this section LAST even though it appears first.
\end{abstract}

% --- Keywords ---
\vspace{0.5em}
\noindent\textbf{Keywords:} % keyword1, keyword2, keyword3

\section{Introduction}
% What is the problem? Why does it matter? What has been done before?
% What is the gap in existing work? What did you do to address it?

\section{Related Work}
% Survey of prior work. Use real citations only. BibTeX format.
% If you cannot find a DOI, note the source honestly.

\section{Methods}
% What data did you use? Where did it come from?
% How did you process it? What experiments did you run?
% Be specific enough that someone could reproduce your work.

\section{Results}
% Present your findings using figures and tables.
% Reference every figure in the text. Let the data speak.
%
% Example figure inclusion:
% \begin{figure}[H]
%   \centering
%   \includegraphics[width=\columnwidth]{figures/figure-1.png}
%   \caption{Description of what this figure shows.}
%   \label{fig:result1}
% \end{figure}

\section{Discussion}
% What do the results mean? What are the limitations?
% What surprised you? What would you do differently?
% Be honest about weaknesses.

\section{Conclusion}
% Brief. What did you find? What is the takeaway?
% What should someone investigate next?

\section*{Acknowledgments}
% Optional. Credit data sources, tools, or collaborators.

\bibliography{references}

\end{document}
